\section{Introduction (not done yet)} \label{sec:Intro}

Holotomography is a problem naturally arising at synchroton facilities.
When we want to look into very small objects ranging in the nanometer regime, a well understood approach is X-ray near-field holography, which yields significant magnification by the chosen cone-beam geometry.
By rotating the sample, we can reconstruct the three-dimensional structure similar to other tomographic imaging techniques.
But there is a big difference compared to these, as we couple the two ill-posed inverse problems phase retrieval and tomography.
\\
This poses computational challenges, because the rays going through the object interact in free space, but also from a theoretical point of view the problem becomes harder as the ambiguity of phase retrieval adds to the typical challenges of tomography.

When working on inverse problems, one goal is always to reach a high resolution in your reconstructed objects.
But as this comes at cubic cost for 3D objects, it is desirable to represent the data on a more flexible structure with strong expressivity using less parameters.
Furthermore, the number of parameters should not increase polynomially with the resolution, as most objects have backgrounds or regions without any structure that do not benefit from higher resolution.
The idea of using neural networks to represent the object, known as \textit{implicit neural representation} (INR) was applied earlier to image compression \cite{donyNeuralNetworkApproaches1995}, 3D scenes \cite{rezendeUnsupervisedLearning3D2018, sitzmannImplicitNeuralRepresentations2020, mullerInstantNeuralGraphics2022} and inverse problems as computed tomography \cite{zhaNAFNeuralAttenuation2022, zhangDynamicConebeamCT2023, essakineWhereWeStand2025}.\\
INRs come at the advantage that the representation is not constrained to a certain grid and can be queried at arbitrary spatial coordinates to obtain the corresponding value of the object.
A crucial point is of course the choice of the network architecture.
As e.g. shown in \cite[Fig. 2(a)]{mullerInstantNeuralGraphics2022}, the Lipschitz continuity of multilayer perceptrons (MLPs) with rectified linear unit (ReLU) activation functions does not allow for sharp edges which are substantial in many images and physical objects.
This tendency to learn low frequencies better is commonly referred to as \textit{spectral bias} \cite{rahamanSpectralBiasNeural2019}.
One way to overcome the inability to represent sharp edges was to focus on activation functions such as \( x \mapsto \sin(x) \) \cite{sitzmannImplicitNeuralRepresentations2020} or similar for which there is an extensive overview by Essakine \textit{et al.} \cite{essakineWhereWeStand2025}. 
Another approach to omit  is to encode the coordinates before feeding them into an MLP with ReLU activation functions.
This was done by Mildenhall \textit{et al.} via a \textit{frequency encoding} \cite{mildenhallNeRFRepresentingScenes2020}.
While the representations with this architecture are powerful and could also be used in holotomography the training process takes a lot of time, which makes them less suited than the architecture chosen by Müller \textit{et al.} we will apply in this paper \cite{mullerInstantNeuralGraphics2022}. 
